\chapter{GARCH e Volatilidade Condicional}
\label{cap:garch}
% ── índice ──────────────────────────────────────────────────────
\index{GARCH|textbf}
\index{volatilidade condicional|textbf}
\index{heterocedasticidade condicional|textbf}
\index{ARCH|see{GARCH}}
\index{garch@\texttt{garch()}|textbf}
\index{agrupamento de volatilidade}
\index{EGARCH}
\index{GJR-GARCH}

% ─────────────────────────────────────────────────────────────────────────────
\section{Heterocedasticidade condicional}

Séries financeiras exibem um fenômeno recorrente: períodos de alta
volatilidade tendem a ser seguidos de mais volatilidade, e períodos calmos
persistem. Esse \emph{agrupamento de volatilidade} (\textit{volatility
clustering}) viola a hipótese de variância condicional constante dos modelos
ARMA — a variância de $\varepsilon_t$ dado o passado varia ao longo do tempo.

O modelo GARCH(\textit{p,q}) (Engle 1982; Bollerslev 1986) captura essa
dinâmica explicitamente. A especificação GARCH(1,1) é:

\[
  \varepsilon_t \mid \mathcal{F}_{t-1} \sim (0,\, h_t)
  \qquad
  h_t = \omega + \alpha_1\,\varepsilon_{t-1}^2 + \beta_1\,h_{t-1}
\]

\noindent onde $h_t$ é a \emph{variância condicional}, $\alpha_1$ captura o
efeito ARCH (reação ao choque passado) e $\beta_1$ captura a persistência
(memória da volatilidade).

\begin{center}
\begin{tabular}{lll}
\toprule
\textbf{Restrição} & \textbf{Condição} & \textbf{Significado} \\
\midrule
Variância positiva  & $\omega > 0,\; \alpha_1 \geq 0,\; \beta_1 \geq 0$ & $h_t > 0$ garantido \\
Estacionariedade    & $\alpha_1 + \beta_1 < 1$                           & variância não explode \\
Variância incondicional & $\sigma^2 = \omega/(1-\alpha_1-\beta_1)$       & nível de longo prazo \\
\bottomrule
\end{tabular}
\end{center}

O caso $\alpha_1 + \beta_1 = 1$ é chamado IGARCH (\emph{integrated}): a
volatilidade tem raiz unitária e choques têm efeito permanente.

% ─────────────────────────────────────────────────────────────────────────────
\section{Teste ARCH-LM}

Antes de estimar um modelo de volatilidade, testa-se se existem efeitos ARCH.
O teste de Engle (1982) regride $\varepsilon_t^2$ em seus próprios lags e
avalia a significância conjunta:

\[
  \varepsilon_t^2 = a_0 + a_1\,\varepsilon_{t-1}^2 + \cdots + a_m\,\varepsilon_{t-m}^2 + u_t
\]

\[
  H_0{:}\; a_1 = \cdots = a_m = 0
  \quad \text{(sem efeitos ARCH)}
\]

A estatística é $LM = N \cdot R^2_{\text{aux}} \xrightarrow{d} \chi^2(m)$
sob $H_0$.

% ─────────────────────────────────────────────────────────────────────────────
\section{Verificação por DGP}

O DGP gera um processo ARCH(1) verdadeiro ($\beta=0$):

\[
  h_t = 0{,}3 + 0{,}5\,\varepsilon_{t-1}^2
  \qquad
  \varepsilon_t = \sqrt{h_t}\; z_t
  \qquad
  z_t \sim (0,1)
\]

A inovação é gerada diretamente como sorteio normal padrão:
$z = \text{rnormal}()$, portanto $E[z]=0$ e $E[z^2]=1$.

\begin{lstlisting}[caption={DGP ARCH(1) e estimação GARCH}]
seed(42)
input df
id
1
end
for i in 1..=9 { df = append(df, df) }
df |> filter(_n <= 500)
generate df z  = rnormal()               // z: E[z]=0, E[z²]=1
generate df e  = z
generate df t  = _n
let n = nrow(df)
let i = 2
while i <= n {
    let le = mean(df, e, if = t == i - 1)
    let lz = mean(df, z, if = t == i)
    replace df e = sqrt(0.3 + 0.5 * le * le) * lz if t == i
    i = i + 1
}

archtest(df, e, lags=5)      // deve rejeitar H₀
let m = garch(df, e, p=1, q=1)
display m
archtest(m, lags=5)          // deve não rejeitar (GARCH captou o ARCH)
\end{lstlisting}

\subsection*{Teste ARCH-LM — série bruta}

A série bruta deve rejeitar a hipótese nula de ausência de ARCH, confirmando
variância condicional variável no processo simulado.

\subsection*{Estimação GARCH(1,1)}

O estimador MLE recupera o DGP verdadeiro ($\omega=0{,}3$, $\alpha=0{,}5$,
$\beta=0$) até a variação de amostra finita. Um $\hat\beta$ próximo de zero
identifica que o processo é ARCH puro — sem componente GARCH de persistência.\footnote{%
  Com $N=500$, estimativas em amostras finitas não precisam coincidir
  exatamente com os parâmetros do DGP; com $N$ maior, a convergência melhora.}

\subsection*{ARCH-LM nos resíduos padronizados}

Após ajustar o GARCH, os resíduos padronizados $\hat{z}_t = \varepsilon_t / \sqrt{\hat{h}_t}$
devem ser ruído branco em $z_t^2$:

Depois do ajuste, o teste ARCH-LM nos resíduos padronizados não deve rejeitar
se a especificação de volatilidade captou a heterocedasticidade condicional.

% ─────────────────────────────────────────────────────────────────────────────
\section{Variantes: EGARCH e GJR-GARCH}

O GARCH simétrico trata choques positivos e negativos igualmente. Em
ativos financeiros, choques negativos tendem a aumentar mais a volatilidade
do que positivos de mesma magnitude (\emph{efeito de alavancagem}).

O \textbf{EGARCH} (Nelson 1991) modela $\ln h_t$:
\[
  \ln h_t = \omega + \beta_1 \ln h_{t-1}
             + \alpha_1 \left|\frac{\varepsilon_{t-1}}{\sqrt{h_{t-1}}}\right|
             + \gamma_1 \frac{\varepsilon_{t-1}}{\sqrt{h_{t-1}}}
\]
O parâmetro $\gamma_1 < 0$ captura a assimetria (efeito de alavancagem).

O \textbf{GJR-GARCH} (Glosten-Jagannathan-Runkle 1993) adiciona um termo
indicador:
\[
  h_t = \omega + \alpha_1 \varepsilon_{t-1}^2
        + \gamma_1 \varepsilon_{t-1}^2 \mathbf{1}[\varepsilon_{t-1} < 0]
        + \beta_1 h_{t-1}
\]
O coeficiente $\gamma_1 > 0$ implica que choques negativos elevam mais a
volatilidade.

% ─────────────────────────────────────────────────────────────────────────────
\section{Sintaxe}

\begin{lstlisting}
garch(df, y, p=1, q=1)
garch(df, y, p=1, q=1, dist=t)    // erros Student-t
egarch(df, y, p=1, q=1)
gjrgarch(df, y, p=1, q=1)
archtest(df, y, lags=5)            // ARCH-LM na série bruta
archtest(modelo, lags=5)           // ARCH-LM nos resíduos padronizados
\end{lstlisting}

\begin{center}
\begin{tabular}{lll}
\toprule
\textbf{Parâmetro} & \textbf{Padrão} & \textbf{Descrição} \\
\midrule
\texttt{p}    & 1                & ordem ARCH ($\alpha$) \\
\texttt{q}    & 1                & ordem GARCH ($\beta$) \\
\texttt{dist} & \texttt{normal}  & distribuição: \texttt{normal} ou \texttt{t} \\
\texttt{lags} & 5                & lags para ARCH-LM \\
\bottomrule
\end{tabular}
\end{center}

\begin{lstlisting}[caption={Fluxo típico — retornos de ativos}]
load "retornos.csv" as df

// 1. Verificar agrupamento de volatilidade
archtest(df, ret, lags=10)

// 2. Estimar modelo de volatilidade
let m_garch  = garch(df, ret, p=1, q=1)
let m_gjr    = gjrgarch(df, ret, p=1, q=1)

// 3. Verificar diagnóstico
archtest(m_garch, lags=10)

// 4. Variância condicional estimada
display m_garch
\end{lstlisting}

\begin{nota}
  O \hayashi{} estima GARCH por MLE assumindo inovações gaussianas
  (\lstinline{dist=normal}) ou Student-t (\lstinline{dist=t}). Para séries
  com caudas muito pesadas (criptomoedas, câmbio em crises), a distribuição
  $t$ geralmente melhora o ajuste. Comparar AIC/BIC entre especificações.
\end{nota}

\section{Validação empírica}

O DGP ARCH(1) deste capítulo corresponde ao caso de validação
\texttt{garch\_book} em \texttt{validation/cases/}. O caso compara os
coeficientes MLE do \hayashi{} com a referência Python \texttt{arch}. Execute
\lstinline{hay validate} para ver o estado atual.
